\section{Conclusion}
\label{sec:conclusion}

The financial accounts of the United States have described the claim structure as a matrix of
counterparties for seventy years. They do not record which income services each claim, and that
omission hides a position nobody chose to take.

Decomposed by servicing income, about \result{DebtOnlyDirect} percent of all US debt, public and
private, is paid directly out of wages, and \result{DebtOnlyIndirect} percent once indirect
servicing is counted. Crossed with the holder structure, the same decomposition shows the federal
government standing behind about \result{SovereignUnion} percent of the directly wage-backed
claims as holder, guarantor or debtor, against about \result{AIFederalShareRounded} percent of
the separately measured claims on AI capital. Counting only what the public sector holds
understates that position by more than a factor of two, which is why the headline carries its
convention wherever it appears.

The long series puts the present level in its place. The federal share traces a U, from
\result{SovUnion1952} percent in 1952 through \result{SovUnion1970} percent in
1970 to \result{SovUnion2025} percent in 2025, so today is a return rather than a
record, and a large part of the recent rise is growth in federal debt itself rather than new
exposure to households. The ratio underneath it is close to flat across seventy-three years. The
share of the national debt stock that wages service is not a recent development.

Across the wage distribution, claims are spread far more evenly than the income that services
them, and the thinnest liquid buffers sit one quintile above the bottom rather than at it. A
proportional fall in pay is therefore regressive before any policy response is chosen, and relief
designed around the bottom of the distribution is aimed one quintile away from the least room.

What the accounts do not do is predict. Two pre-registered tests, the second on
\result{PredPanelRows} bank-years, found that adding the measure to a conventional bank-risk
model makes out-of-sample ranking worse rather than better. We report both because they bound the
contribution: this is accounting that says where wage dependence sits, not a measure that says
which institution will take a loss. Every number here is generated from an artifact in a
replication package, each carrying its status, so a reader who rejects a convention can change it
and see exactly what moves.
